\section{แนวคิดการออกแบบคลังข้อมูล}
\subsection{การแยกงานวิเคราะห์จากงานประจำวัน}
ฐานข้อมูลต้นทางเก็บรายการและข้อมูลประกอบตามเหตุการณ์ ส่วนคลังข้อมูลจัดเก็บข้อมูลที่ผ่านการแปลงและเชื่อมรหัสอ้างอิงเพื่ออ่านซ้ำในหลายรายงาน เส้นทางข้อมูลประกอบด้วยการสกัด การแปลง และการโหลด หรือ ETL จากนั้นบริการรายงานจำกัดขอบเขตข้อมูลก่อนจัดรูปแบบสำหรับหน้าเว็บ

\subsection{ข้อเท็จจริง มิติ และระดับรายละเอียด}
ข้อเท็จจริงเก็บเหตุการณ์และค่าที่วัดได้ ส่วนมิติเก็บคุณลักษณะสำหรับจัดกลุ่มและกรองข้อมูล การกำหนดระดับรายละเอียดหรือ grain ต้องมาก่อนการเลือกกราฟ เพราะใบลาเดียวอาจสัมพันธ์กับหลายรายวิชา หลายกลุ่มเรียน และหลายผู้สอน การเชื่อมทุกความสัมพันธ์แล้วรวมวันลาทันทีทำให้ยอดภาพรวมสูงเกินจริง

แบบจำลองรุ่นใหม่จึงมีลักษณะกลุ่มดาวข้อเท็จจริง โดยใช้หลายตารางข้อเท็จจริงร่วมมิติเดียวกัน แยกตารางเชื่อมสำหรับความสัมพันธ์หลายต่อหลาย และแยกความเป็นเจ้าของรายวิชาออกจากสังกัดนิสิตและบุคลากร เพื่อไม่ให้รายงานต่างระดับใช้คำว่า ``ภาควิชา'' แทนข้อมูลคนละความหมาย

\subsection{บริบทในอดีตและสิทธิ์ปัจจุบัน}
สังกัด หลักสูตร และที่ปรึกษา ณ เวลายื่นเป็นบริบทเชิงวิเคราะห์ ขณะที่รายชื่อผู้ที่ผู้ใช้เข้าถึงได้ต้องตรวจจากสิทธิ์และความสัมพันธ์ปัจจุบัน การมีรหัสบุคคลอยู่ในคลังข้อมูลไม่ได้ให้สิทธิ์เปิดดูข้อมูลของบุคคลนั้นโดยอัตโนมัติ
